\begin{lemma}[Shared realization of a finite randomized learner]\label{lem:step-011-shared-randomization}
For every randomized map
\(A:(X_{k,N}\times\{0,1\})^n\to\mathcal H_{k,N}\), there is a
deterministic function

\[
  \Phi_A:(X_{k,N}\times\{0,1\})^n\times(0,1)
  \longrightarrow\mathcal H_{k,N}
\]

such that, for \(V\sim\operatorname{Unif}(0,1)\),
\(\Phi_A(s,V)\) has exactly the law \(A(s)\) for every
\(s\in(X_{k,N}\times\{0,1\})^n\). For the same \(V\), equal input datasets give equal
outputs pathwise. No properness, determinism, or tag symmetry of \(A\) is
required.
\end{lemma}

\begin{proof}
Enumerate the finite set
\(\mathcal H_{k,N}=\{h_1,\ldots,h_q\}\). For each dataset \(s\), put

\[
  c_m(s):=\sum_{a=1}^m\Pr\{A(s)=h_a\},\qquad c_0(s):=0,
\]

and let \(\Phi_A(s,v)=h_m\) at the first index satisfying
\(c_m(s)\ge v\). Since \(c_q(s)=1\), this is total, and

\[
  \Pr\{\Phi_A(s,V)=h_m\}=c_m(s)-c_{m-1}(s)
  =\Pr\{A(s)=h_m\}.
\]

This proves the exact learner marginal. Equal datasets and the same \(V\)
give the same two arguments to the same deterministic function, proving
the pathwise synchronization statement. \end{proof}

\begin{lemma}[Fixed-hidden common-randomness coupling]\label{lem:step-011-coupling-kernel}
Under Propositions~\ref{prop:step-009-overflow} and
\ref{prop:step-010-simulator} and
Lemma~\ref{lem:step-011-shared-randomization}, fix any \((t,Q)\) in (2)
and any arbitrary randomized \(A\). There exists a joint probability
kernel carrying

\begin{enumerate}
\item An infinite sequence
   \[
     W_\ell=(X_\ell^\star,Y_\ell^\star),\qquad \ell\ge1,
   \tag{7}
   \]
   of mutually independent draws from \(Q^{\tau_t}\), so
   \(X_\ell^\star\sim Q\) and
   \(Y_\ell^\star=\tau_t(X_\ell^\star)\);
\item Independently of (7), the complete input-independent preprocessing seed
   of Lemma~\ref{lem:step-010-public-preprocessing}; and
\item Independently of both, one \(V\sim\operatorname{Unif}(0,1)\).
\end{enumerate}

such that the following statements hold.

\begin{itemize}
\item The actual input is
  \[
    Z^{\mathrm{in}}_\ell=W_\ell,\qquad \ell=1,\ldots,M,
  \tag{8}
  \]
  and hence has exactly the law (3), independently of simulator
  preprocessing.
\item Put \(\Theta_J=(T_J,Q_J):=(t,Q)\), retain the seed's
  \(\Theta_i=(T_i,Q_i)\) for \(i\ne J\), write
  \(\boldsymbol\Xi=(\Theta_1,\ldots,\Theta_k)\), and define the always-present
  ideal ordered dataset
  \[
    S^{\mathrm{id}}=(\widehat z_1,\ldots,\widehat z_n)
  \tag{9}
  \]
  by
  \[
    \widehat z_r=
    \begin{cases}
      ((J,X_{L_r}^\star),Y_{L_r}^\star),&I_r=J,\\[1mm]
      ((I_r,X_r^\circ),Y_r^\circ),&I_r\ne J,
    \end{cases}
    \qquad r\in[n].
  \tag{10}
  \]
  The infinite sequence makes (10) well defined even when \(U>M\).
\item There is a deterministic realization \(\Phi_A\) of the learner kernel
  for which
  \[
    H^{\mathrm{id}}:=\Phi_A(S^{\mathrm{id}},V),
    \qquad
    G^{\mathrm{id}}:=D_JH^{\mathrm{id}}
  \tag{11}
  \]
  has the exact ideal call-and-restrict marginal.
\item Define the actual output by the simulator control flow:
  \[
    G^{\mathrm{act}}:=
    \begin{cases}
      g_0,&U>M,\\
      D_J\Phi_A(\mathcal P_{\Omega_{\mathrm{pre}}}
                    (Z^{\mathrm{in}}),V),&U\le M.
    \end{cases}
  \tag{12}
  \]
  Its marginal is exactly
  \(B_{\mu_{N,M},A}(Z^{\mathrm{in}})\) with input law (3).
  In the first line of (12), \(V\), (7), and every other coupling-only
  variable are ignored: operationally the actual simulator still reads no
  input and makes no call to \(A\).
\end{itemize}

Thus (7)-(12) give a coupling of the exact actual and ideal experiments,
with every expectation and source of randomness specified.
\end{lemma}

\begin{proof}
The finite law \(Q^{\tau_t}\) defines the product sequence (7). Only its
first \(M\) coordinates are exposed as the actual input, so (8) is exactly
an iid size-\(M\) input rather than a stopped or conditioned sample. By
construction this sequence is independent of the preprocessing seed, as
required by Lemma~\ref{lem:step-010-public-preprocessing}.

The running indices satisfy \(1\le L_r\le U\le n\) at every hidden
position. Hence the infinite sequence gives every hidden row in (10),
including on overflow. Nonhidden rows in (10) are the identical realized
rows already present in the preprocessing seed. Therefore (9) is
always an ordered member of the full size-\(n\) learner domain.

Lemma~\ref{lem:step-011-shared-randomization} makes (11) an exact call to
the randomized kernel \(A\), followed by
the same restriction used by the simulator. In (12), the
nonoverflow branch applies the same realization to the exact row
map, while the overflow branch is the exact fixed output from Proposition~\ref{prop:step-010-simulator}. Sampling an unused coupling
variable does not create an operational learner call. Equations (8) and
(12), together with the input-independent seed law and that lemma, prove the
actual marginal claim. Equations (9)-(11) define the ideal marginal.
\end{proof}

\begin{lemma}[Identification with the ideal experiment]\label{lem:step-011-ideal-marginal}
Under Lemmas~\ref{lem:step-008-instance-factorization} and
\ref{lem:step-008-ideal-sample-law}, and under
Lemma~\ref{lem:step-011-coupling-kernel}, fix \(j\in[k]\) and a complete
vector

\[
  \boldsymbol\theta=((t_1,Q_1),\ldots,(t_k,Q_k))
\tag{14}
\]

with \((t_j,Q_j)=(t,Q)\). Conditional on \(J=j\) and on the nonhidden
seed instances taking the values in (14), but without conditioning on
\(U\), \(U\le M\), or any tag-count event, the ideal dataset (9) satisfies

\[
  S^{\mathrm{id}}\sim
  \left(P_{\boldsymbol Q}^{c_{\boldsymbol t}}\right)^n.
\tag{15}
\]

Equivalently, the \(n\) coupled ideal global labeled rows are mutually independent,
and each has the fixed realizable law that first selects a tag uniformly
and then draws its feature from that tag's fixed \(Q_i\). If the nonhidden
instances are not fixed, their marginal is only the corresponding mixture
of the product laws in (15); this lemma does not call that mixture iid from
one deterministic product instance. Moreover, if the hidden pair
\((t,Q)\) is first drawn from \(\mu_{N,M}\), the complete joint marginal
\((J,\boldsymbol\Xi,S^{\mathrm{id}},H^{\mathrm{id}})\) is exactly the
ideal experiment of Section~\ref{app:step-008}, and

\[
  \mathbb E_{(t,Q)\sim\mu_{N,M}}
    \mathbb E_{\mathrm{id},(t,Q)}R_Q(G^{\mathrm{id}},\tau_t)
  =\mathbb E_{\mathrm{ideal}}
    R_{Q_J}(D_JH^{\mathrm{id}},\tau_{T_J}).
\tag{15a}
\]
\end{lemma}

\begin{proof}
Fix ordered tags \(i_1,\ldots,i_n\in[k]\) and features
\(x_1,\ldots,x_n\in[N]\). The tag kernel gives

\[
  \Pr\{I_1=i_1,\ldots,I_n=i_n\mid J=j\}=k^{-n}.
\tag{16}
\]

Conditional on this tag vector and (14), enumerate the positions with
tag \(j\) in increasing order. Formula (10) assigns them the distinct
sequence variables \(X_1^\star,\ldots,X_u^\star\), where \(u\) is the
number of such positions. Their joint mass is the product of the
corresponding \(Q_j\) factors. At every position with tag \(i_r\ne j\),
Lemma~\ref{lem:step-010-public-preprocessing} supplies an
independent feature with law \(Q_{i_r}\). Those nonhidden features are
conditionally mutually independent and are independent of the separately
drawn hidden sequence (7). Therefore

\[
  \Pr\{\text{the ideal feature at row }r\text{ is }x_r
        \text{ for all }r
       \mid J=j,\boldsymbol\theta,I_{1:n}=i_{1:n}\}
  =\prod_{r=1}^n Q_{i_r}(x_r).
\tag{17}
\]

All labels in (10) are deterministic and equal
\(\tau_{t_{i_r}}(x_r)\). Multiplying (16) and (17) gives

\[
  \prod_{r=1}^n\left[
    \frac1k Q_{i_r}(x_r)
    \mathbf1\{y_r=\tau_{t_{i_r}}(x_r)\}
  \right]
\tag{18}
\]

for every ordered labeled row vector. Expression (18) is exactly the
product mass function in (15). This calculation sums over all tag vectors,
including those with \(U>M\); no overflow conditioning appears.

Now draw the hidden pair \((t,Q)\sim\mu_{N,M}\) independently of the
selector before applying the fixed-hidden kernel. Conditional on \(J=j\),
this pair occupies coordinate \(j\), while
Lemma~\ref{lem:step-010-public-preprocessing} draws every coordinate
\(i\ne j\) independently from the
same \(\mu_{N,M}\). These are precisely the hypotheses of Lemma~\ref{lem:step-008-instance-factorization}; hence the full vector has
law \(\mu_{N,M}^{\otimes k}\) and is independent of \(J\). The row kernel
verified in (16)-(18) is precisely the kernel in Lemma~\ref{lem:step-008-ideal-sample-law}. Finally,
Lemma~\ref{lem:step-011-shared-randomization} gives the exact same learner
kernel \(A\). Thus the complete coupled ideal marginal equals the ideal experiment, and its selected restriction is \(D_JH^{\mathrm{id}}\),
which proves (15a). Lemma~\ref{lem:step-008-selector-independence} and
Proposition~\ref{prop:step-008-selected-risk-identity} therefore remain
valid for this marginal; their exchangeability and product-risk
conclusions are consumed as facts and are not rederived here.

Conditioning (15) further on \(U\le M\) would restrict the tag vectors in
(16) and generally destroy the iid tag law. The proof never makes that
conditional-iid assertion: later equality on \(U\le M\) is pathwise under
the coupling, while (15) is the unconditioned ideal sample law for the fixed
block vector. \end{proof}

\begin{proposition}[Exact row and restricted-output identity off overflow]\label{prop:step-011-no-overflow-identity}
Under Lemma~\ref{lem:step-011-coupling-kernel}, Proposition~\ref{prop:step-010-row-construction}, and Lemma~\ref{lem:step-010-one-use}, on the event

\[
  E:=\{U\le M\},
\tag{19}
\]

the actual ordered dataset and ideal ordered dataset satisfy, row by row,

\[
  \mathcal P_{\Omega_{\mathrm{pre}}}(Z^{\mathrm{in}})_r
  =\widehat z_r,
  \qquad r=1,\ldots,n.
\tag{20}
\]

Consequently the shared randomized learner outputs and restrictions in
(11)-(12) agree pathwise on \(E\), and

\[
  R_Q(G^{\mathrm{id}},\tau_t)
  =R_Q(G^{\mathrm{act}},\tau_t)
  \qquad\text{on }E.
\tag{21}
\]

This includes \(U=0\) and \(U=M\). No dataset, output, or risk equality is
claimed on \(E^c=\{U>M\}\).
\end{proposition}

\begin{proof}
Fix an outcome in \(E\) and a row \(r\in[n]\). If \(I_r=J\), then

\[
  1\le L_r\le U\le M.
\tag{22}
\]

The exact occurrence map (6) uses actual input row \(L_r\).
By (8), this row is

\[
  Z^{\mathrm{in}}_{L_r}
  =(X_{L_r}^\star,Y_{L_r}^\star).
\]

The first line of (6) is therefore

\[
  ((J,X_{L_r}^\star),Y_{L_r}^\star),
\]

which is exactly the first line of the ideal rule (10). If \(I_r\ne J\),
the second lines of (6) and (10) are literally the same realized
nonhidden row \(((I_r,X_r^\circ),Y_r^\circ)\). This exhaustive row split
proves (20), with ordering preserved.

At \(U=0\), there is no hidden row and the second-line identity applies at
all \(n\) positions; both experiments call \(A\) on the identical all-
nonhidden dataset. At \(U=M\), (22) still holds, and the exact occurrence
map may use all and only the first \(M\) sequence records. Thus neither
endpoint requires a separate conditioning argument or an extra record.

Because (20) gives equal datasets and both calls use the same \(V\),
Lemma~\ref{lem:step-011-shared-randomization} gives equal hypotheses even when
\(A\) is randomized, improper, nonmonotone, or tag-asymmetric. Both sides
then apply the same restriction \(D_J\), so the one-block outputs are
equal. Evaluating the identical functions against the same fixed
\((Q,\tau_t)\) proves (21).

On \(E^c\), Proposition~\ref{prop:step-010-simulator} makes the
actual side return \(g_0\) before reading the input or calling \(A\),
whereas the ideal side still uses (9)-(11). The proof makes no comparison
between those datasets or outputs there. \end{proof}

\begin{proposition}[Bounded selected-risk transfer]\label{prop:step-011-risk-transfer}
Under Lemma~\ref{lem:step-011-ideal-marginal},
Proposition~\ref{prop:step-011-no-overflow-identity}, and Proposition~\ref{prop:step-009-overflow}, for every fixed hidden pair
\((t,Q)\) and arbitrary randomized \(A\), define

\[
  L_{\mathrm{id}}:=R_Q(G^{\mathrm{id}},\tau_t),
  \qquad
  L_{\mathrm{act}}:=R_Q(G^{\mathrm{act}},\tau_t).
\tag{23}
\]

Then, on the full coupling space,

\[
  L_{\mathrm{id}}
  \ge L_{\mathrm{act}}-\mathbf1\{U>M\}.
\tag{24}
\]

Taking the exact marginals from
Lemma~\ref{lem:step-011-coupling-kernel} proves (4a). Averaging the hidden
pair under \(\mu_{N,M}\) and applying
Lemma~\ref{lem:step-011-ideal-marginal} proves the exact ideal
interface (4b). Proposition~\ref{prop:step-009-overflow} supplies
the strict residual (5).
If \(k=2\) or \(k=3\), then \(U\le M\) pathwise and the same result has
zero residual.
\end{proposition}

\begin{proof}
Population 0-1 risk is a probability, so for every output function,

\[
  0\le L_{\mathrm{id}}\le1,
  \qquad
  0\le L_{\mathrm{act}}\le1.
\tag{25}
\]

On \(E\), Proposition~\ref{prop:step-011-no-overflow-identity} gives
\(L_{\mathrm{id}}=L_{\mathrm{act}}\), so (24) holds with a zero
indicator. On \(E^c\), (25) gives

\[
  L_{\mathrm{act}}-1\le0\le L_{\mathrm{id}},
\]

which is exactly (24) with indicator one. This overflow argument uses only
boundedness. In particular, the actual early-abort output is \(g_0\), but
its risk may be any value in \([0,1]\): it is zero for the all-zero target,
one for the all-one target under any \(Q\), and can take an intermediate
value for other thresholds. No equality or favorable abort-risk property
is assumed.

Taking expectation of (24) over the joint variables in
Lemma~\ref{lem:step-011-coupling-kernel} gives

\[
  \mathbb E_{\mathrm{id},\theta}L_{\mathrm{id}}
  \ge
  \mathbb E_{\mathrm{act},\theta}L_{\mathrm{act}}
  -\mathbb E\mathbf1\{U>M\}.
\tag{26}
\]

The actual marginal identity in that lemma identifies the middle term with
the simulator expectation displayed in (4a), including all of its
preprocessing and learner randomness. The ideal marginal identifies the
first term with the ideal selected-risk expectation. Finally,

\[
  \mathbb E\mathbf1\{U>M\}=\Pr\{U>M\}<2^{-9}
\tag{27}
\]

by Proposition~\ref{prop:step-009-overflow}. There is no conditioning-to-expectation
conversion and no interchange of a limit, supremum, or integral: the
ideal dataset uses only finitely many coordinates of (7), and all other
spaces are finite.

Because (26) holds for every fixed \((t,Q)\), average it over
\((t,Q)\sim\mu_{N,M}\). The overflow term is unchanged because its law
depends only on \(J,I_{1:n}\). Lemma~\ref{lem:step-011-ideal-marginal},
specifically (15a), identifies the averaged ideal term with the exact
selected-risk expectation in Section~\ref{app:step-008}. The averaged
actual term is
the right-hand simulator expectation in (4b). This proves (4b) without
rederiving or using the equality between ideal selected risk and
product population risk.

For \(k=2,3\), Proposition~\ref{prop:step-009-overflow} gives
\(M\ge n\ge U\), so (21) holds everywhere and the final term in (26) is
exactly zero. For \(n<k\) or \(M=8\), the coupling and pointwise inequality
are unchanged, while the floor-eight estimate supplies the stated
strict residual. \end{proof}
